\chapter*{\textbf{LISTA DE ABREVIAÇÕES E ACRÔNIMOS}}
\begin{acronym}[ROCOF] %% Depois de rodar a primeira vez, digite entre os [ ] a sigla mais longa para justificar a lista de acrônimos à esquerda
\acro{AM}{Amplitude Modulation}
\acro{AR}{Autoregressive}
\acro{BS}{B-Spline}
\acro{DFT}{Discrete Fourier Transform}
\acro{FFT}{Fast Fourier Transform}
\acro{FIR}{Finite Impulse Response}
%\acro{Fs}{Sampling Frequency}
%\acro{f0}{Nominal Frequency}
%\acro{fm}{Modulation Frequency}
%\acro{H}{Number of Harmonics}
\acro{IEC}{International Electrotechnical Commission}
\acro{IEEE}{Institute of Electrical and Electronics Engineers}
%\acro{ka}{Phase Modulation Index}
%\acro{kx}{Amplitude Modulation Index}
%\acro{Nc}{Number of Cycles}
%\acro{Nppc}{Number of Points per Cycle}
\acro{PM}{Phase Modulation}
\acro{PMU}{Phasor Measurement Unit}
\acro{ROCOF}{Rate of Change of Frequency}
\acro{RMS}{Root Mean Square}
\acro{SNR}{Signal-to-Noise Ratio}
\acro{TFT}{Taylor-Fourier Transform}
\acro{TVE}{Total Vector Error}
\acro{WAMS}{Wide-Area Measurement Systems}
\acro{ZC}{Zero Crossing}
\acro{ZCD}{Zero Crossing Detector}

\end{acronym}